\chapter{Conclusions and Future Work}

\section{Conclusions}

This thesis demonstrates a variable-height UAV CKM predictor under a strict
city-holdout protocol. The final system jointly predicts dense path loss, delay
spread, and angular spread for unseen cities while conditioning on continuous
UAV height. The strongest result is not a pure black-box network, but a hybrid
pipeline: calibrated priors explain the stable physical and statistical
structure, and a neural residual model corrects what remains.

The final numbers meet the original targets. The path-loss prior reaches
\SI{1.9246}{dB} RMSE on its evaluation holdout and \SI{1.9383}{dB} when used as
the frozen prior on the final test split. The spread priors reach
\SI{28.10}{ns} DS RMSE and \(13.74^\circ\) AS RMSE. The final residual GMM-head
model improves those anchors on 14 unseen test cities, reaching
\SI{1.6519}{dB}, \SI{26.5570}{ns}, and \(11.3854^\circ\) for PL, DS, and AS.
The prior is therefore not a weak baseline; it is already a strong estimator.

The answer to \textbf{RQ1} is yes within the fixed CKM contract: dense PL/DS/AS
prediction is viable under strict city holdout, and the final model does it
jointly with one generator. The answer to \textbf{RQ2} is that the useful split
is calibrated physics/statistics first and bounded neural residuals second. The
priors combine two-ray LoS behaviour, height-aware calibration, NLoS morphology,
and log-domain DS/AS ridge models; the network learns local residuals around
these anchors. The answer to \textbf{RQ3} is also positive: continuous UAV
height is handled through height-aware priors and a sinusoidal FiLM conditioning
path, while low-altitude and NLoS-heavy cases remain the most difficult.

The main contributions are therefore:
\begin{itemize}
  \item a joint CKM generator for variable-height UAV path loss, delay spread,
        and angular spread on unseen cities;
  \item a calibrated prior family that combines physical propagation structure,
        morphology, visibility, height regimes, and empirical spread statistics;
  \item a distribution-first modelling lesson showing why GMM-style residual
        prediction and histogram checks are needed for NLoS path-loss modes and
        spread tails;
  \item a reproducible city-holdout evaluation contract with ground-only
        masking and LoS/NLoS/height diagnostics.
\end{itemize}

The qualitative spread results also clarify the remaining limitation. The
general DS/AS structure is usually correct, but rare groups of high-spread
pixels are sometimes missed, and some false positive high-spread patches remain.
This explains why spread RMSE improves less dramatically than path loss even
though the maps are physically more coherent.

The negative results justify the final design. Pure image-translation models,
larger CNN variants, topology experts, multi-scale outputs, and uncertainty
heads exposed useful failure modes, but none solved the distribution mismatch
alone. The final model is stronger because it combines auditable priors with a
distribution-aware residual head instead of only increasing neural capacity.

\section{Future directions}

The nearest continuation is a controlled ablation of the final residual line:
remove individual prior channels, residual caps, GMM terms, anchor/guard losses,
and height-conditioning paths while keeping the same city-holdout contract.
This would quantify which parts of the final gain are essential.

The next modelling step is broader conditioning. The current dataset fixes
frequency, antenna assumptions, and simulator settings; a more general CKM
surrogate should learn across frequency, bandwidth, antenna pattern, and
possibly polarization. A practical intermediate step would simulate the same
cities under several radio configurations so the model can separate geometry
effects from radio-configuration effects.

For spread targets, future work should also test structure-aware objectives.
DS and AS maps are quantized and spike-plus-tail, so small spatial shifts can be
penalized harshly by RMSE even when the morphology is plausible. Region-wise
correlation, gradient consistency, LoS/NLoS boundary consistency, or
topology-aware ranking losses may better preserve physical spread patterns.

The final extension is simulator-to-reality calibration. Because the priors are
interpretable, they are suitable for correction with limited field measurements.
A future operational CKM predictor should adapt to a new city with sparse real
data while still reporting uncertainty by LoS/NLoS regime, UAV-height bin, and
topology class.
